%%%%%%%%%%%%%%%%%%%%%%%%%%%%%%%%%%%%%%%%%%%%%%%%%%%%%%%%%%%%
%% Relaxed Concurrent Counting Bloom Filter — FPR & FNR  %%
%% Full Proofs                                            %%
%%%%%%%%%%%%%%%%%%%%%%%%%%%%%%%%%%%%%%%%%%%%%%%%%%%%%%%%%%%%
\documentclass[11pt,a4paper]{article}

%% ── Packages ────────────────────────────────────────────
\usepackage{amsmath,amssymb,amsthm}
\usepackage{hyperref}

%% ── Theorem environments ────────────────────────────────
\newtheorem{theorem}{Theorem}[section]
\newtheorem{lemma}[theorem]{Lemma}
\newtheorem{corollary}[theorem]{Corollary}
\newtheorem{proposition}[theorem]{Proposition}
\newtheorem{definition}[theorem]{Definition}
\newtheorem{remark}[theorem]{Remark}

%% ── Notation shortcuts ─────────────────────────────────
\newcommand{\CBF}{\textsf{CBF}}
\newcommand{\RCCBF}{\textsf{RCCBF}}
\newcommand{\FPR}{\mathrm{FPR}}
\newcommand{\FNR}{\mathrm{FNR}}
\newcommand{\Poi}{\mathrm{Poi}}
\newcommand{\E}{\mathbb{E}}
\newcommand{\Prob}{\mathbb{P}}
\newcommand{\R}{\mathbb{R}}
\newcommand{\N}{\mathbb{N}}
\newcommand{\calB}{\mathcal{B}}

%% ── Title ───────────────────────────────────────────────
\title{%
  False Positive and False Negative Rate Analysis\\
  for the Relaxed Concurrent Counting Bloom Filter}
\author{%
  % TODO: Replace with actual author names and affiliations
  Author One\thanks{Department of Computer Science, University X.}
  \and
  Author Two\thanks{Department of Computer Science, University Y.}
}
\date{March 2026}

\begin{document}
\maketitle

%%%%%%%%%%%%%%%%%%%%%%%%%%%%%%%%%%%%%%%%%%%%%%%%%%%%%%%%%%%%
%% ABSTRACT
%%%%%%%%%%%%%%%%%%%%%%%%%%%%%%%%%%%%%%%%%%%%%%%%%%%%%%%%%%%%
\begin{abstract}
We present a complete and self-contained analysis of the false positive
rate~(FPR) and false negative rate~(FNR) for the \emph{Relaxed Concurrent
Counting Bloom Filter}~(\RCCBF{}), a probabilistic data structure that
replaces each counter in a standard counting Bloom filter with a
\emph{MultiCounter} composed of $p$ sub-counters updated via the
power-of-two-choices paradigm.  We derive asymptotically exact expressions
in the mean-field limit ($p\to\infty$) for the FPR and
FNR under the Poisson hashing model, establish closed-form series for the
threshold-one case using the mean-field ODE of Mitzenmacher~\cite{mitzenmacher1996},
prove that the \RCCBF{} strictly improves the FPR over a standard counting
Bloom filter for every relaxation parameter $p>1$, and characterise the
FNR--FPR trade-off as a function of the load and the relaxation parameter.
Special cases recovering the classical Bloom filter behaviour are discussed.
\end{abstract}

%%%%%%%%%%%%%%%%%%%%%%%%%%%%%%%%%%%%%%%%%%%%%%%%%%%%%%%%%%%%
%% SECTION 1 — INTRODUCTION
%%%%%%%%%%%%%%%%%%%%%%%%%%%%%%%%%%%%%%%%%%%%%%%%%%%%%%%%%%%%
\section{Introduction}\label{sec:intro}

Bloom filters and their counting variants are fundamental building blocks
in high-throughput systems such as network routers, distributed caches,
and database engines, where membership queries must be answered in
constant time with bounded memory.  As data rates grow and multi-core
architectures become prevalent, concurrent access to these structures is
no longer optional but essential.

\paragraph{The problem.}
Standard counting Bloom filters~(\CBF{}) maintain an array of $m$
counters; each insertion increments $k$ counters, and each query checks
whether those counters all meet a threshold.  Under concurrent access,
each counter update requires an atomic compare-and-swap (CAS) operation.
When the load is high and many threads target the same counter, CAS
contention degrades throughput and introduces long-tail latencies.

\paragraph{The idea.}
We study the \emph{Relaxed Concurrent Counting Bloom Filter}~(\RCCBF{}),
which replaces each counter with a \emph{MultiCounter}---an array of $p$
sub-counters updated via the power-of-two-choices paradigm.  An increment
selects two sub-counters uniformly at random and increments the lighter
one; a read samples one sub-counter uniformly and scales by~$p$.  This
distributes write contention across $p$ locations while still providing
a useful estimate of the total count.  The relaxation parameter~$p$
controls the trade-off between concurrency and accuracy.

\paragraph{Contributions.}
Our main contributions are as follows.
\begin{itemize}
  \item \textbf{FPR formula.}  We derive an asymptotically exact
        expression for the false positive rate of the \RCCBF{} in the
        mean-field limit ($p \to \infty$) and prove that it strictly
        improves upon the standard Bloom filter FPR for all $p > 1$
        (Theorems~\ref{thm:fpr-rccbf} and~\ref{thm:fpr-strict}).
  \item \textbf{FNR formula.}  We derive a corresponding expression for
        the false negative rate, which is zero in a standard \CBF{} but
        non-zero in the \RCCBF{} due to the stochastic read operation
        (Theorem~\ref{thm:fnr-rccbf}).
  \item \textbf{Regime analysis.}  We characterise the FPR--FNR trade-off
        across sparse, crossover, and dense regimes, identifying the
        operating points where the \RCCBF{} offers the greatest advantage
        (Propositions~\ref{prop:sparse}--\ref{prop:dense} and
        Table~\ref{tab:regimes}).
\end{itemize}

\paragraph{Paper outline.}
Section~\ref{sec:prelims} introduces notation and the Poisson hashing
model.  Section~\ref{sec:cbf} recapitulates the classical Bloom filter
FPR.  Section~\ref{sec:multicounter} develops the MultiCounter and its
mean-field analysis.  Section~\ref{sec:construction} defines the \RCCBF{}
formally.  Section~\ref{sec:fpr} derives the FPR.
Section~\ref{sec:fnr} derives the FNR.  Section~\ref{sec:special}
verifies special cases and summarises the parameter regimes.
Section~\ref{sec:conclusion} concludes.

%%%%%%%%%%%%%%%%%%%%%%%%%%%%%%%%%%%%%%%%%%%%%%%%%%%%%%%%%%%%
%% SECTION 2 — PRELIMINARIES
%%%%%%%%%%%%%%%%%%%%%%%%%%%%%%%%%%%%%%%%%%%%%%%%%%%%%%%%%%%%
\section{Preliminaries and Notation}\label{sec:prelims}

\begin{definition}[Counting Bloom Filter Parameters]\label{def:params}
A \emph{counting Bloom filter}~$\calB$ is parameterised by:
\begin{itemize}
  \item $m \in \N$: the number of counter positions,
  \item $n \in \N$: the number of elements inserted,
  \item $k \in \N$: the number of independent hash functions
        $h_1, h_2, \ldots, h_k$, each mapping the universe to
        $\{1,2,\ldots,m\}$ uniformly and independently.
\end{itemize}
\end{definition}

\begin{definition}[Relaxation Parameters]\label{def:relaxation}
A \emph{Relaxed Concurrent Counting Bloom Filter}~($\RCCBF$) extends the
counting Bloom filter with:
\begin{itemize}
  \item $p \in \N$, $p \ge 1$: the number of \emph{sub-counters} per
        position (the \emph{relaxation parameter}),
  \item $\theta \in \N$: the \emph{query threshold}—a query returns
        ``present'' only if the estimated count is at least~$\theta$,
  \item $\tau = \lceil \theta / p \rceil$: the \emph{effective sub-counter
        threshold}.
\end{itemize}
\end{definition}

\begin{definition}[Expected Load]\label{def:load}
The \emph{expected load} per position is
\begin{equation}\label{eq:lambda}
  \lambda \;=\; \frac{kn}{m}.
\end{equation}
Under the Poisson hashing model (i.e.\ each element independently hashes to
each position with probability $1/m$, so the count at any position is
approximately $\Poi(\lambda)$ for large $m$), this is the parameter
governing the occupancy distribution at each counter.
\end{definition}

\begin{remark}[Hash Function Assumptions]\label{rem:hash}
Throughout, we assume the hash functions $h_1, \ldots, h_k$ are
\emph{fully independent} and \emph{uniform} over $\{1,\ldots,m\}$.
Under this assumption, the values $h_i(x)$ and $h_j(y)$ are independent
for all distinct elements~$x \ne y$ and all indices~$i, j$.
\end{remark}

%%%%%%%%%%%%%%%%%%%%%%%%%%%%%%%%%%%%%%%%%%%%%%%%%%%%%%%%%%%%
%% SECTION 3 — STANDARD CBF FPR
%%%%%%%%%%%%%%%%%%%%%%%%%%%%%%%%%%%%%%%%%%%%%%%%%%%%%%%%%%%%
\section{Standard Counting Bloom Filter}\label{sec:cbf}

We briefly recapitulate the classical false positive analysis of the
(standard, non-counting) Bloom filter, following Bloom~\cite{bloom1970}.
The counting variant, introduced by Fan et al.~\cite{fan2000}, replaces
bits with counters to support deletions; its false positive behaviour is
identical to the standard Bloom filter when the query threshold is~$1$.

\begin{definition}[Standard Bloom Filter Operations]\label{def:bf-ops}
In a standard Bloom filter with bit array of size $m$:
\begin{itemize}
  \item \textbf{Insert}$(x)$: For $i = 1, \ldots, k$, set bit $h_i(x)$
        to~$1$.
  \item \textbf{Query}$(y)$: Return ``present'' if and only if bit $h_i(y)
        = 1$ for all $i = 1, \ldots, k$.
\end{itemize}
A counting Bloom filter~\cite{fan2000} replaces bits with counters,
incrementing on insert and decrementing on delete, and queries whether the
counter is at least~$1$.
\end{definition}

\begin{theorem}[Classical Bloom Filter FPR]\label{thm:bf-fpr}
Under the Poisson hashing model, the false positive rate of a standard
Bloom filter is
\begin{equation}\label{eq:fpr-bf}
  \FPR_{\mathrm{BF}}
  \;=\; \bigl(1 - e^{-\lambda}\bigr)^{k},
  \qquad \lambda = \frac{kn}{m}.
\end{equation}
\end{theorem}

\begin{proof}
Consider an element $y$ \emph{not} in the stored set~$S$.
For a single position~$j \in \{1,\ldots,m\}$, the number of elements in~$S$
whose hash maps to position~$j$ is, under the Poisson model,
$C_j \sim \Poi(\lambda)$.  The probability that position~$j$ is \emph{non-zero}
is
\begin{equation}\label{eq:pj-nonzero}
  \Prob(C_j \ge 1) \;=\; 1 - \Prob(C_j = 0)
  \;=\; 1 - e^{-\lambda}.
\end{equation}
Because the $k$ hash functions are independent, the $k$ positions
$h_1(y), \ldots, h_k(y)$ are independent.  A false positive occurs when
all $k$ positions are non-zero:
\[
  \FPR_{\mathrm{BF}}
  = \prod_{i=1}^{k}\Prob(C_{h_i(y)} \ge 1)
  = \bigl(1 - e^{-\lambda}\bigr)^{k}.
\]
\end{proof}

\begin{remark}[Optimal Number of Hash Functions]\label{rem:optimal-k}
The FPR is minimised when the derivative with respect to $k$ vanishes.
Setting $f(k) = (1 - e^{-kn/m})^k$ and solving
$\frac{d}{dk}\ln f(k) = 0$ yields the well-known optimum
\begin{equation}\label{eq:optimal-k}
  k^{*} = \frac{m}{n}\ln 2.
\end{equation}
At this optimum, $\FPR_{\mathrm{BF}}^{*} = (1/2)^k = 2^{-k}$.
\end{remark}

%%%%%%%%%%%%%%%%%%%%%%%%%%%%%%%%%%%%%%%%%%%%%%%%%%%%%%%%%%%%
%% SECTION 4 — THE MULTICOUNTER
%%%%%%%%%%%%%%%%%%%%%%%%%%%%%%%%%%%%%%%%%%%%%%%%%%%%%%%%%%%%
\section{The MultiCounter (Power-of-Two-Choices)}\label{sec:multicounter}

The key building block of the \RCCBF{} is the \emph{MultiCounter}: an
array of $p$ sub-counters updated using the balanced-allocation
(power-of-two-choices) strategy introduced by Azar et
al.~\cite{azar1999} and analysed via mean-field ODEs by
Mitzenmacher~\cite{mitzenmacher1996,mitzenmacher2001}.

\begin{definition}[MultiCounter Operations]\label{def:multicounter}
A MultiCounter $M = (S_1, S_2, \ldots, S_p)$ supports:
\begin{itemize}
  \item \textbf{Increment}: Choose two sub-counters $S_a, S_b$ uniformly
        at random (with replacement).  Increment whichever of $S_a$, $S_b$
        has the smaller current value (ties broken arbitrarily).
  \item \textbf{Read}: Choose one sub-counter $S_j$ uniformly at random.
        Return the \emph{scaled value} $R = p \cdot S_j$.
\end{itemize}
\end{definition}

\begin{remark}\label{rem:total-count}
If $C$ total increments have been applied to the MultiCounter, then
$\sum_{j=1}^{p} S_j = C$.  The mean sub-counter value is $C/p$.
The read operation returns an \emph{estimate} of the total count scaled to
compensate for sampling.
\end{remark}

\subsection{Mean-Field Analysis}\label{ssec:mean-field}

We now derive the occupancy tail probabilities of a MultiCounter via
the mean-field ODE framework.  This follows Mitzenmacher~\cite{mitzenmacher1996}
and rests on the convergence result of Kurtz~\cite{kurtz1970} for
density-dependent Markov chains.

\begin{definition}[Occupancy Tail Fractions]\label{def:tail}
For a MultiCounter with $p$ sub-counters and total count $C$, define
\begin{equation}\label{eq:si-def}
  s_i(\mu) \;=\; \lim_{p \to \infty}
  \frac{1}{p}\,\bigl|\{j : S_j \ge i\}\bigr|,
  \qquad \mu = \frac{C}{p},
\end{equation}
with the convention $s_0(\mu) = 1$ for all $\mu \ge 0$.
Intuitively, $s_i(\mu)$ is the fraction of sub-counters whose value is at
least~$i$, as a function of the average load~$\mu$.
\end{definition}

\begin{theorem}[Mean-Field ODE, Mitzenmacher~\cite{mitzenmacher1996}]
\label{thm:mean-field-ode}
As $p \to \infty$, the occupancy tail fractions satisfy the system of
differential equations
\begin{equation}\label{eq:ode-system}
  \frac{ds_i}{d\mu}
  \;=\; s_{i-1}(\mu)^2 - s_i(\mu)^2,
  \qquad i = 1, 2, 3, \ldots,
\end{equation}
with initial conditions $s_i(0) = 0$ for all $i \ge 1$ and
boundary condition $s_0(\mu) = 1$ for all $\mu \ge 0$.
\end{theorem}

The convergence of the finite-$p$ system to the ODE~\eqref{eq:ode-system}
follows from Kurtz's theorem~\cite{kurtz1970} on density-dependent Markov
chains; see also the survey by Mitzenmacher, Richa, and
Sitaraman~\cite{mitzenmacher2001survey}.

\begin{remark}[Finite-$p$ Convergence Rate]\label{rem:finite-p-rate}
By Kurtz's theorem~\cite{kurtz1970}, the convergence rate of the finite-$p$
system to the mean-field ODE~\eqref{eq:ode-system} is $O(1/\sqrt{p})$ for
sample paths and $O(1/p)$ for expectations.  For $p \ge 8$, the relative
error in $s_1$ is below $1\%$.  Table~\ref{tab:finite-p} illustrates this
for $C = 1$.
\end{remark}

\begin{table}[h]
\centering
\begin{tabular}{ccccc}
\hline
$p$ & $C$ & Exact $\Prob(S_j \ge 1)$ & $\tanh(C/p)$ & Relative Error \\
\hline
2  & 1 & 0.500  & 0.462 & 7.6\% \\
4  & 1 & 0.250  & 0.245 & 2.0\% \\
8  & 1 & 0.125  & 0.124 & 0.5\% \\
16 & 1 & 0.0625 & 0.0624 & 0.14\% \\
\hline
\end{tabular}
\caption{Finite-$p$ accuracy of the mean-field approximation $s_1(C/p)
  = \tanh(C/p)$ for the probability that a uniformly random sub-counter is
  at least~$1$, with $C = 1$ total increment.}
\label{tab:finite-p}
\end{table}

\subsection{Explicit Solution for \texorpdfstring{$i=1$}{i=1}}
\label{ssec:tanh}

\begin{lemma}\label{lem:s1-tanh}
The solution to the ODE~\eqref{eq:ode-system} for $i = 1$ is
\begin{equation}\label{eq:s1-tanh}
  s_1(\mu) \;=\; \tanh(\mu).
\end{equation}
\end{lemma}

\begin{proof}
For $i = 1$, the ODE~\eqref{eq:ode-system} with $s_0 = 1$ becomes
\begin{equation}\label{eq:ode-i1}
  \frac{ds_1}{d\mu} = 1 - s_1^2,
  \qquad s_1(0) = 0.
\end{equation}
This is a separable ODE.  Separating variables:
\begin{equation}\label{eq:sep}
  \frac{ds_1}{1 - s_1^2} = d\mu.
\end{equation}
We use the partial-fraction decomposition
\begin{equation}\label{eq:partial-frac}
  \frac{1}{1 - s_1^2}
  = \frac{1}{(1 - s_1)(1 + s_1)}
  = \frac{1}{2}\left(\frac{1}{1 - s_1} + \frac{1}{1 + s_1}\right).
\end{equation}
Integrating both sides:
\begin{align}
  \int \frac{ds_1}{1 - s_1^2}
  &= \frac{1}{2}\int\frac{ds_1}{1 - s_1}
     + \frac{1}{2}\int\frac{ds_1}{1 + s_1} \notag \\
  &= -\frac{1}{2}\ln|1 - s_1|
     + \frac{1}{2}\ln|1 + s_1| + \kappa \notag \\
  &= \frac{1}{2}\ln\frac{1 + s_1}{1 - s_1} + \kappa.
  \label{eq:int-lhs}
\end{align}
Here $\kappa$ denotes the integration constant.  Setting this equal to
$\mu + \kappa$ and applying the initial condition $s_1(0) = 0$:
\begin{equation}\label{eq:ic}
  \frac{1}{2}\ln\frac{1 + 0}{1 - 0} = 0 + \kappa
  \implies \kappa = 0.
\end{equation}
Therefore
\begin{equation}\label{eq:solve-for-s1}
  \frac{1}{2}\ln\frac{1 + s_1}{1 - s_1} = \mu
  \implies \frac{1 + s_1}{1 - s_1} = e^{2\mu}.
\end{equation}
Solving for $s_1$:
\begin{align}
  1 + s_1 &= e^{2\mu}(1 - s_1), \notag \\
  1 + s_1 &= e^{2\mu} - e^{2\mu} s_1, \notag \\
  s_1(1 + e^{2\mu}) &= e^{2\mu} - 1, \notag \\
  s_1 &= \frac{e^{2\mu} - 1}{e^{2\mu} + 1}
       = \tanh(\mu). \label{eq:s1-final}
\end{align}
The last equality follows from the identity
$\tanh(\mu) = (e^{2\mu}-1)/(e^{2\mu}+1)$.
\end{proof}

\subsection{Approximations for General \texorpdfstring{$\tau$}{tau}}
\label{ssec:general-tau}

\begin{lemma}[Small-Load Approximation]\label{lem:small-load}
For $\mu \ll 1$ and integer $\tau \ge 1$,
\begin{equation}\label{eq:small-load}
  s_\tau(\mu)
  \;\approx\;
  \frac{\mu^{2^\tau - 1}}{c_\tau},
\end{equation}
where $c_\tau$ is a constant depending only on $\tau$, defined by the
recursion $c_\tau = (2^\tau - 1)\,c_{\tau-1}^2$ with $c_1 = 1$.
In particular:
\begin{equation}\label{eq:c-tau-values}
  c_1 = 1, \qquad c_2 = 3, \qquad c_3 = 63, \qquad c_4 = 59535.
\end{equation}
\end{lemma}

\begin{proof}
We proceed by induction on $\tau$.  For $\tau = 1$, we have
$s_1(\mu) = \tanh(\mu)$.  For small $\mu$, $\tanh(\mu) = \mu - \mu^3/3
+ O(\mu^5)$, so $s_1(\mu) \approx \mu = \mu^{2^1 - 1}/c_1$ with
$c_1 = 1$.

For the inductive step, assume $s_{\tau-1}(\mu) \approx
\mu^{2^{\tau-1}-1}/c_{\tau-1}$ for small $\mu$.
From the ODE~\eqref{eq:ode-system}:
\begin{equation}\label{eq:ode-inductive}
  \frac{ds_\tau}{d\mu}
  = s_{\tau-1}^2 - s_\tau^2
  \approx s_{\tau-1}^2
  \approx \frac{\mu^{2(2^{\tau-1}-1)}}{c_{\tau-1}^2}
  = \frac{\mu^{2^{\tau}-2}}{c_{\tau-1}^2},
\end{equation}
where we dropped the $s_\tau^2$ term because $s_\tau \ll s_{\tau-1}$ in
the small-$\mu$ regime.  Integrating with $s_\tau(0) = 0$:
\begin{equation}\label{eq:s-tau-approx}
  s_\tau(\mu)
  \approx \frac{\mu^{2^\tau - 1}}{(2^\tau - 1)\,c_{\tau-1}^2}.
\end{equation}
Setting $c_\tau = (2^\tau - 1)\,c_{\tau-1}^2$ gives the claimed form.
The explicit values follow by unwinding the recursion:
$c_2 = 3 \cdot 1^2 = 3$,
$c_3 = 7 \cdot 3^2 = 63$,
$c_4 = 15 \cdot 63^2 = 59535$.
\end{proof}

\begin{lemma}[Large-Load Concentration]\label{lem:large-load}
When $\mu = C/p \gg 1$, all sub-counters are concentrated in the interval
\begin{equation}\label{eq:concentration}
  S_j \;\in\; \Bigl[\frac{C}{p} - O(\log\log p),\;
  \frac{C}{p} + O(\log\log p)\Bigr]
\end{equation}
with probability $1 - o(1)$ as $p \to \infty$.
\end{lemma}

\begin{proof}
This is the celebrated result of Azar et al.~\cite{azar1999} for balanced
allocations with two choices: the maximum load of any bin exceeds the
average by at most $O(\log\log p)$ with high probability.
In the heavily loaded case ($C \gg p$), Berenbrink et
al.~\cite{berenbrink2006} sharpen this to show that the gap
between maximum and minimum load is $O(\log\log p)$ with high probability.
Therefore every sub-counter satisfies~\eqref{eq:concentration}, and the
read of a uniformly random sub-counter returns a value within
$O(p \cdot \log\log p)$ of the true total count~$C$.
See also Peres, Talwar, and Wieder~\cite{peres2015} for refined bounds
in the graphical balanced-allocation setting.
\end{proof}

%%%%%%%%%%%%%%%%%%%%%%%%%%%%%%%%%%%%%%%%%%%%%%%%%%%%%%%%%%%%
%% SECTION 5 — RCCBF CONSTRUCTION
%%%%%%%%%%%%%%%%%%%%%%%%%%%%%%%%%%%%%%%%%%%%%%%%%%%%%%%%%%%%
\section{RCCBF --- Construction}\label{sec:construction}

We now formally define the Relaxed Concurrent Counting Bloom Filter.

\begin{remark}[Space Complexity]\label{rem:space}
The \RCCBF{} uses $m \cdot p$ sub-counters.  If each sub-counter uses $b$
bits, the total memory is
\begin{equation}\label{eq:space}
  M_{\mathrm{bits}} = m \cdot p \cdot b \text{ bits},
\end{equation}
compared with $m \cdot b$ for a standard \CBF{}---a $p\times$ overhead.
\end{remark}

\begin{definition}[\RCCBF{}]\label{def:rccbf}
An \RCCBF{} is a tuple $(\calB, m, k, p, \theta)$ where:
\begin{itemize}
  \item $\calB = (M_1, M_2, \ldots, M_m)$ is an array of $m$
        MultiCounters, each containing $p$ sub-counters
        (Definition~\ref{def:multicounter}).
  \item $h_1, \ldots, h_k : \mathcal{U} \to \{1, \ldots, m\}$ are $k$
        fully independent, uniform hash functions.
  \item $\theta \ge 1$ is the query threshold.
\end{itemize}

The operations are defined as follows:
\begin{itemize}
  \item \textbf{Insert}$(x)$: For $i = 1, \ldots, k$, invoke
        $M_{h_i(x)}.\mathrm{Increment}()$.
  \item \textbf{Query}$(y)$: For $i = 1, \ldots, k$, let
        $R_i = M_{h_i(y)}.\mathrm{Read}()$.
        Return ``present'' if and only if $R_i \ge \theta$ for all
        $i = 1, \ldots, k$.
\end{itemize}
\end{definition}

\begin{remark}[Concurrency]\label{rem:concurrency}
The MultiCounter increment uses only an atomic compare-and-swap on the
chosen sub-counter, making Insert lock-free.  The Read operation is a
single atomic load, hence wait-free.  This yields a data structure that is
\emph{distributionally linearizable} in the sense of Alistarh et
al.~\cite{alistarh2018}: individual operations are not linearizable, but
the \emph{distribution} of outcomes matches a sequential specification up
to bounded relaxation controlled by~$p$.
\end{remark}

\begin{remark}[Concurrent Execution and Mean-Field Validity]
\label{rem:concurrent-validity}
The analysis in this paper assumes sequential allocation via the mean-field
ODE (Theorem~\ref{thm:mean-field-ode}).  Under concurrent execution, stale
reads during the two-choice selection may degrade load balance: a thread
may observe an outdated sub-counter value and increment a sub-counter that
is no longer the lightest.  Alistarh et al.~\cite{alistarh2018} establish
an $O(n \log^2 n)$ distributional linearizability bound for relaxed
concurrent structures of this kind.  A full concurrent analysis
incorporating stale-read effects on the occupancy distribution is deferred
to future work.
\end{remark}

\begin{remark}[Reduction to Standard CBF]\label{rem:p1}
When $p = 1$ and $\theta = 1$, the MultiCounter degenerates to a single
counter.  Increment adds~1 to the counter, and Read returns the counter
value.  The \RCCBF{} then coincides with a standard counting Bloom
filter.
\end{remark}

%%%%%%%%%%%%%%%%%%%%%%%%%%%%%%%%%%%%%%%%%%%%%%%%%%%%%%%%%%%%
%% SECTION 6 — FPR ANALYSIS
%%%%%%%%%%%%%%%%%%%%%%%%%%%%%%%%%%%%%%%%%%%%%%%%%%%%%%%%%%%%
\section{False Positive Rate Analysis}\label{sec:fpr}

This section contains the main contribution: a complete derivation
of the false positive rate of the \RCCBF{}.

\subsection{Setup and Statement}\label{ssec:fpr-setup}

Let $S = \{x_1, \ldots, x_n\}$ be the set of inserted elements, and let
$y \notin S$ be an element that was \emph{never} inserted.  A
\emph{false positive} occurs when $\mathrm{Query}(y)$ returns ``present''.

\begin{theorem}[RCCBF False Positive Rate]\label{thm:fpr-rccbf}
In the mean-field limit ($p \to \infty$), under the Poisson hashing model
with parameters $m, n, k, p, \theta$ and $\lambda = kn/m$, the false
positive probability for $y \notin S$ is
\begin{equation}\label{eq:fpr-rccbf}
  \boxed{%
  \FPR_{\RCCBF}
  \;=\; q(\lambda, p, \tau)^{k},}
\end{equation}
where $\tau = \lceil \theta / p \rceil$ and
\begin{equation}\label{eq:q-def}
  q(\lambda, p, \tau)
  \;=\; \sum_{c=0}^{\infty}
        \frac{e^{-\lambda}\,\lambda^c}{c!}\;\cdot\;
        s_\tau\!\left(\frac{c}{p}\right).
\end{equation}
\end{theorem}

The remainder of this section is devoted to the proof.

\subsection{Step 1: Poisson Count at a Non-Member Position}
\label{ssec:fpr-step1}

\begin{lemma}\label{lem:poisson-count}
For any position $j \in \{1, \ldots, m\}$ and any element $y \notin S$,
the total count $C_j$ at position~$j$ due to the $n$ inserted elements is
\begin{equation}\label{eq:poisson-count}
  C_j \;\sim\; \Poi(\lambda), \qquad \lambda = \frac{kn}{m}.
\end{equation}
\end{lemma}

\begin{proof}
Each inserted element $x_\ell \in S$ contributes to position $j$ through
hash function $h_i$ independently with probability $1/m$.  Across all $n$
elements and all $k$ hash functions, there are $kn$ independent Bernoulli
trials, each with success probability $1/m$.  The total count at
position~$j$ is therefore $C_j \sim \mathrm{Bin}(kn, 1/m)$.

Under the Poisson approximation (valid for large $m$ with $kn/m = \lambda$
held constant, and exact in the Poisson hashing model),
\begin{equation}\label{eq:bin-to-poi}
  \mathrm{Bin}(kn, 1/m) \;\to\; \Poi(\lambda),
\end{equation}
which is the standard balls-into-bins limit.

Crucially, the element~$y$ was \emph{never} inserted, so $y$ contributes
zero to every position.  The count $C_j$ arises entirely from hash
collisions with the $n$~elements of~$S$.
\end{proof}

\subsection{Step 2: Independence Across Positions}
\label{ssec:fpr-step2}

\begin{lemma}\label{lem:independence}
Under fully independent hash functions, the counts
$C_{h_1(y)}, C_{h_2(y)}, \ldots, C_{h_k(y)}$ are independent random
variables, each distributed as $\Poi(\lambda)$.
\end{lemma}

\begin{proof}
The positions $h_1(y), \ldots, h_k(y)$ are independent uniform random
variables over $\{1, \ldots, m\}$ (by the independence of the hash
functions).  Under the Poisson hashing model, the counts at distinct
positions are independent Poisson random variables.  Even when two
positions coincide (which happens with probability $O(1/m)$), the Poisson
model treats each hash function's contribution independently.  Therefore
$C_{h_1(y)}, \ldots, C_{h_k(y)}$ are i.i.d.\ $\Poi(\lambda)$.
\end{proof}

\begin{remark}[Hash Collision Correction]\label{rem:hash-collision}
The probability that any two of the $k$ positions $h_1(y), \ldots, h_k(y)$
coincide is $O(k^2/m)$.  We suppress this correction throughout; for
typical parameter choices ($k \ll \sqrt{m}$), the error is negligible.
\end{remark}

\subsection{Step 3: Read Exceeding Threshold Given Count}
\label{ssec:fpr-step3}

\begin{lemma}\label{lem:read-threshold}
Let $C$ increments be applied to a MultiCounter with $p$ sub-counters
via the two-choice protocol.  Let $R = p \cdot S_j$ be the read value,
where $S_j$ is a uniformly random sub-counter.  Then, in the mean-field
limit ($p \to \infty$),
\begin{equation}\label{eq:read-prob}
  \Prob(R \ge \theta \mid C)
  \;=\; \Prob\!\bigl(S_j \ge \tau \mid C\bigr)
  \;=\; s_\tau\!\left(\frac{C}{p}\right),
\end{equation}
where $\tau = \lceil \theta / p \rceil$.
\end{lemma}

\begin{proof}
The read operation selects a sub-counter $S_j$ uniformly at random from
$\{S_1, \ldots, S_p\}$.  The query threshold is met when $R = p S_j \ge
\theta$, i.e.\ $S_j \ge \theta / p$, which (since $S_j$ is integer-valued)
is equivalent to $S_j \ge \lceil \theta / p \rceil = \tau$.

By Definition~\ref{def:tail}, the fraction of sub-counters with value at
least~$\tau$, in the mean-field limit, is $s_\tau(C/p)$.  Since $S_j$ is
chosen uniformly:
\begin{equation}\label{eq:uniform-choice}
  \Prob(S_j \ge \tau \mid C)
  = \frac{|\{j : S_j \ge \tau\}|}{p}
  \;\xrightarrow{p \to \infty}\; s_\tau\!\left(\frac{C}{p}\right).
\end{equation}
\end{proof}

\subsection{Step 4: Unconditioning Over the Poisson Count}
\label{ssec:fpr-step4}

\begin{lemma}\label{lem:uncondition}
The probability that the read at a single position of $y \notin S$ exceeds
the threshold is
\begin{equation}\label{eq:q-uncondition}
  q(\lambda, p, \tau)
  \;=\; \E\!\left[s_\tau\!\left(\frac{C}{p}\right)\right]
  \;=\; \sum_{c=0}^{\infty}
        \frac{e^{-\lambda}\,\lambda^c}{c!}\;
        s_\tau\!\left(\frac{c}{p}\right),
\end{equation}
where $C \sim \Poi(\lambda)$.
\end{lemma}

\begin{proof}
By the law of total probability, conditioning on the Poisson count $C$ at
the position and applying Lemma~\ref{lem:read-threshold}:
\begin{align}
  q(\lambda, p, \tau)
  &= \Prob(R \ge \theta) \notag \\
  &= \sum_{c=0}^{\infty} \Prob(R \ge \theta \mid C = c)\;\Prob(C = c)
     \notag \\
  &= \sum_{c=0}^{\infty} s_\tau\!\left(\frac{c}{p}\right)\;
     \frac{e^{-\lambda}\,\lambda^c}{c!}.
  \label{eq:q-expansion}
\end{align}
Note that $s_\tau(0) = 0$ for all $\tau \ge 1$, so the $c = 0$ term
vanishes.
\end{proof}

\subsection{Step 5: Product Over \texorpdfstring{$k$}{k} Positions}
\label{ssec:fpr-step5}

\begin{proof}[Proof of Theorem~\ref{thm:fpr-rccbf}]
A false positive occurs when the read at \emph{every} one of the $k$
positions of $y$ meets or exceeds the threshold.  By
Lemma~\ref{lem:independence}, the counts at the $k$ positions are
independent.  Conditional on the count at each position, the read
probabilities are determined by Lemma~\ref{lem:read-threshold}, and entries
at different positions are independent (they reside in different
MultiCounters).  Therefore:
\begin{align}
  \FPR_{\RCCBF}
  &= \Prob\!\bigl(R_1 \ge \theta \;\wedge\; R_2 \ge \theta
     \;\wedge\; \cdots \;\wedge\; R_k \ge \theta\bigr) \notag \\
  &= \prod_{i=1}^{k} \Prob(R_i \ge \theta) \notag \\
  &= \prod_{i=1}^{k} q(\lambda, p, \tau) \notag \\
  &= q(\lambda, p, \tau)^k.
  \label{eq:fpr-product}
\end{align}
This completes the proof.
\end{proof}

\subsection{Closed-Form Series for \texorpdfstring{$\tau = 1$}{tau=1}}
\label{ssec:fpr-closedform}

When $\theta \le p$ (so $\tau = 1$), Lemma~\ref{lem:s1-tanh} gives
$s_1(\mu) = \tanh(\mu)$, and we can derive a closed-form series
for~$q$.

\begin{theorem}[Closed-Form FPR for $\tau=1$]\label{thm:fpr-tau1}
For $\tau = 1$,
\begin{equation}\label{eq:q-tau1}
  q(\lambda, p, 1)
  \;=\; (1 - e^{-\lambda})
        - 2\sum_{j=0}^{\infty} (-1)^j\,
        \Bigl[\exp\!\bigl(-\lambda\!\bigl(1 - e^{-2(j+1)/p}\bigr)\bigr)
              - e^{-\lambda}\Bigr].
\end{equation}
\end{theorem}

\begin{proof}
From Lemma~\ref{lem:s1-tanh} and the identity
\begin{equation}\label{eq:tanh-identity}
  \tanh(\mu) = 1 - \frac{2}{e^{2\mu} + 1},
\end{equation}
we write
\begin{equation}\label{eq:q-split}
  q = \E[\tanh(C/p)]
    = \E\!\left[1 - \frac{2}{e^{2C/p} + 1}\right]
    = 1 - 2\,\E\!\left[\frac{1}{e^{2C/p} + 1}\right].
\end{equation}

\medskip\noindent\textbf{Handling $C = 0$ separately.}
Since $s_1(0) = \tanh(0) = 0$, the $c = 0$ term contributes $0$ to $q$.
We split the expectation:
\begin{align}
  \E\!\left[\frac{1}{e^{2C/p} + 1}\right]
  &= \Prob(C = 0) \cdot \frac{1}{e^{0} + 1}
     + \E\!\left[\frac{1}{e^{2C/p} + 1}\,\mathbf{1}_{C \ge 1}\right]
     \notag \\
  &= \frac{e^{-\lambda}}{2}
     + \E\!\left[\frac{1}{e^{2C/p} + 1}\,\mathbf{1}_{C \ge 1}\right].
  \label{eq:split-c0}
\end{align}

\medskip\noindent\textbf{Geometric series for $c \ge 1$.}
For $c \ge 1$, write $u = e^{2c/p} > 1$ (strictly, since $c \ge 1$ and
$p \ge 1$).  Then $|u^{-1}| < 1$ and the geometric series converges
absolutely:
\begin{equation}\label{eq:geometric}
  \frac{1}{u + 1}
  = \frac{u^{-1}}{1 + u^{-1}}
  = \sum_{j=0}^{\infty} (-1)^j\, u^{-(j+1)}
  = \sum_{j=0}^{\infty} (-1)^j\, e^{-2(j+1)c/p}.
\end{equation}

Taking the Poisson expectation over $c \ge 1$ term-by-term (justified by
Fubini's theorem, since the double sum converges absolutely due to geometric decay in~$j$):
\begin{align}
  \E\!\left[\frac{1}{e^{2C/p} + 1}\,\mathbf{1}_{C \ge 1}\right]
  &= \sum_{j=0}^{\infty} (-1)^j
     \sum_{c=1}^{\infty}
     \frac{e^{-\lambda}\,\lambda^c}{c!}\;
     e^{-2(j+1)c/p}.
  \label{eq:swap-sums}
\end{align}
Using the Poisson moment-generating function
$\E[e^{tC}] = \exp(\lambda(e^t - 1))$ for $C \sim \Poi(\lambda)$:
\begin{equation}\label{eq:poi-mgf}
  \sum_{c=0}^{\infty}
  \frac{e^{-\lambda}\,\lambda^c}{c!}\; e^{-2(j+1)c/p}
  = \exp\!\left(-\lambda\!\left(1 - e^{-2(j+1)/p}\right)\right).
\end{equation}
The sum starting at $c = 1$ is obtained by subtracting the $c = 0$ term:
\begin{equation}\label{eq:poi-mgf-sub}
  \sum_{c=1}^{\infty}
  \frac{e^{-\lambda}\,\lambda^c}{c!}\; e^{-2(j+1)c/p}
  = \exp\!\left(-\lambda\!\left(1 - e^{-2(j+1)/p}\right)\right)
    - e^{-\lambda}.
\end{equation}
Substituting~\eqref{eq:poi-mgf-sub} into~\eqref{eq:swap-sums}:
\begin{equation}\label{eq:E-c-ge-1}
  \E\!\left[\frac{1}{e^{2C/p} + 1}\,\mathbf{1}_{C \ge 1}\right]
  = \sum_{j=0}^{\infty} (-1)^j
    \Bigl[\exp\!\bigl(-\lambda\!\bigl(1 - e^{-2(j+1)/p}\bigr)\bigr)
          - e^{-\lambda}\Bigr].
\end{equation}

\medskip\noindent\textbf{Combining.}
From~\eqref{eq:q-split} and~\eqref{eq:split-c0}:
\begin{align}
  q &= 1 - 2\left[\frac{e^{-\lambda}}{2}
       + \sum_{j=0}^{\infty}(-1)^j
         \Bigl(\exp\!\bigl(-\lambda\!\bigl(1 - e^{-2(j+1)/p}\bigr)\bigr)
               - e^{-\lambda}\Bigr)\right] \notag \\
    &= 1 - e^{-\lambda}
       - 2\sum_{j=0}^{\infty}(-1)^j
         \Bigl[\exp\!\bigl(-\lambda\!\bigl(1 - e^{-2(j+1)/p}\bigr)\bigr)
               - e^{-\lambda}\Bigr] \notag \\
    &= (1 - e^{-\lambda})
       - 2\sum_{j=0}^{\infty}(-1)^j
         \Bigl[\exp\!\bigl(-\lambda\!\bigl(1 - e^{-2(j+1)/p}\bigr)\bigr)
               - e^{-\lambda}\Bigr].
  \label{eq:q-final}
\end{align}

\medskip\noindent\textbf{Convergence.}
Each bracketed term tends to $0$ as $j \to \infty$, since
$e^{-2(j+1)/p} \to 0$ implies
$\exp(-\lambda(1 - e^{-2(j+1)/p})) \to e^{-\lambda}$.
Moreover, the bracketed terms are positive, decrease monotonically in $j$
(because $1 - e^{-2(j+1)/p}$ increases in~$j$), and have alternating
signs $(-1)^j$.  By the alternating series test, the series converges.
\end{proof}

\subsection{Strict Improvement Over Standard CBF}\label{ssec:fpr-strict}

\begin{theorem}[Strict FPR Improvement]\label{thm:fpr-strict}
For all $p > 1$ and $\lambda > 0$,
\begin{equation}\label{eq:fpr-strict}
  \FPR_{\RCCBF}(\theta = 1) \;<\; \FPR_{\mathrm{BF}}.
\end{equation}
That is, $q(\lambda, p, 1) < 1 - e^{-\lambda}$.
\end{theorem}

\begin{proof}
Recall $q(\lambda, p, 1) = \E[\tanh(C/p)]$ with $C \sim \Poi(\lambda)$.
The standard CBF false positive probability per position is
$\Prob(C \ge 1) = 1 - e^{-\lambda}$.  We need to show that
$\E[\tanh(C/p)] < 1 - e^{-\lambda}$ for all $p > 1$.

\medskip\noindent\textbf{Step 1.}
We claim that $\tanh(c/p) < \mathbf{1}_{c \ge 1}$ for all integers
$c \ge 0$ and $p > 1$, where $\mathbf{1}_{c \ge 1}$ is the indicator that
$c \ge 1$.

For $c = 0$: $\tanh(0) = 0 = \mathbf{1}_{0 \ge 1}$.  Equality holds.

For $c \ge 1$ and $p > 1$: We have $c/p > 0$ but $c/p$ is finite, so
$\tanh(c/p) < 1 = \mathbf{1}_{c \ge 1}$, since $\tanh(x) < 1$ for all
finite $x > 0$.

\medskip\noindent\textbf{Step 2.}
Taking expectations:
\begin{align}
  q(\lambda, p, 1)
  &= \E[\tanh(C/p)] \notag \\
  &< \E[\mathbf{1}_{C \ge 1}] \notag \\
  &= \Prob(C \ge 1) \notag \\
  &= 1 - e^{-\lambda}.
  \label{eq:strict-ineq}
\end{align}
The strict inequality holds because for $C \ge 1$ (which occurs with
probability $1 - e^{-\lambda} > 0$ for $\lambda > 0$), we have
$\tanh(C/p) < 1 = \mathbf{1}_{C \ge 1}$.

\medskip\noindent\textbf{Step 3.}
Since $q(\lambda, p, 1) < 1 - e^{-\lambda}$ and both sides are positive:
\begin{equation}\label{eq:fpr-strict-final}
  \FPR_{\RCCBF}
  = q(\lambda, p, 1)^k
  < (1 - e^{-\lambda})^k
  = \FPR_{\mathrm{BF}}.
\end{equation}
\end{proof}

\subsection{Scaling Analysis}\label{ssec:fpr-scaling}

\begin{proposition}[Sparse Regime]\label{prop:sparse}
When $\lambda \ll 1$ and $p \gg 1$ (sparse filter with large relaxation
parameter), the FPR ratio satisfies
\begin{equation}\label{eq:sparse-ratio}
  \frac{\FPR_{\RCCBF}}{\FPR_{\mathrm{BF}}}
  \;\approx\; \left(\frac{1}{p}\right)^k.
\end{equation}
\end{proposition}

\begin{proof}
In the sparse regime, $C$ is $0$ or $1$ with high probability.
Specifically, $\Prob(C = 0) = e^{-\lambda} \approx 1 - \lambda$, and
$\Prob(C = 1) \approx \lambda$.  The dominant contribution to $q$ comes
from $c = 1$:
\begin{equation}\label{eq:q-sparse}
  q(\lambda, p, 1)
  \approx \lambda \cdot \tanh(1/p)
  \approx \lambda \cdot \frac{1}{p},
\end{equation}
where we used $\tanh(1/p) \approx 1/p$ for $p \gg 1$.
The standard CBF gives $1 - e^{-\lambda} \approx \lambda$.  Therefore:
\begin{equation}\label{eq:ratio-sparse}
  \frac{q(\lambda, p, 1)}{1 - e^{-\lambda}}
  \approx \frac{\lambda/p}{\lambda} = \frac{1}{p}.
\end{equation}
Raising to the $k$-th power gives the claimed ratio.
\end{proof}

\begin{proposition}[Dense Regime]\label{prop:dense}
When $\lambda \gg p$ (heavily loaded filter),
\begin{equation}\label{eq:dense-ratio}
  \frac{\FPR_{\RCCBF}}{\FPR_{\mathrm{BF}}} \;\to\; 1.
\end{equation}
\end{proposition}

\begin{proof}
When $\lambda \gg p$, the Poisson random variable $C$ concentrates around
$\lambda$, and $C/p \gg 1$ with high probability.  By the concentration
result of Lemma~\ref{lem:large-load}, the sub-counters are tightly
concentrated around $C/p$, and $\tanh(C/p) \to 1$ exponentially fast.
More precisely, $1 - \tanh(C/p) = 2e^{-2C/p}/(1 + e^{-2C/p}) \le
2e^{-2C/p}$.  Therefore $q \to (1 - e^{-\lambda}) - \epsilon$ where $\epsilon \to 0$
exponentially fast, and both the \RCCBF{} and the standard
CBF have FPR approaching~$1$ in the heavily loaded regime.  The ratio
thus tends to~$1$.
\end{proof}

%%%%%%%%%%%%%%%%%%%%%%%%%%%%%%%%%%%%%%%%%%%%%%%%%%%%%%%%%%%%
%% SECTION 7 — FNR ANALYSIS
%%%%%%%%%%%%%%%%%%%%%%%%%%%%%%%%%%%%%%%%%%%%%%%%%%%%%%%%%%%%
\section{False Negative Rate Analysis}\label{sec:fnr}

We now analyse the false negative rate: the probability that an element
$x \in S$ is \emph{not} reported as present by the query.

\subsection{Setup}\label{ssec:fnr-setup}

Let $x \in S$ be an element that was inserted into the \RCCBF{}.
Let $C_i$ denote the total count at position $h_i(x)$ for $i = 1, \ldots, k$.

\begin{remark}\label{rem:fnr-cbf}
In a standard counting Bloom filter ($p = 1$, $\theta = 1$), the counter
at each position of an inserted element is at least~$1$, so the filter
\emph{never} produces a false negative: $\FNR_{\mathrm{BF}} = 0$.
In the \RCCBF{}, however, the Read operation samples a single
sub-counter and scales by~$p$, so it may underestimate the true count,
leading to a non-zero FNR.
\end{remark}

\subsection{Statement}\label{ssec:fnr-statement}

\begin{theorem}[RCCBF False Negative Rate]\label{thm:fnr-rccbf}
In the mean-field limit ($p \to \infty$), for an element $x \in S$ with
counts $C_1, C_2, \ldots, C_k$ at its $k$ hash positions, the false
negative probability is
\begin{equation}\label{eq:fnr-rccbf}
  \boxed{%
  \FNR_{\RCCBF}
  \;=\; 1 - \prod_{i=1}^{k} s_\tau\!\left(\frac{C_i}{p}\right).}
\end{equation}
For $\tau = 1$:
\begin{equation}\label{eq:fnr-tau1}
  \FNR_{\RCCBF}
  \;=\; 1 - \prod_{i=1}^{k} \tanh\!\left(\frac{C_i}{p}\right).
\end{equation}
When all counts are equal ($C_i = C$ for all $i$):
\begin{equation}\label{eq:fnr-equal}
  \FNR_{\RCCBF}
  \;=\; 1 - \tanh\!\left(\frac{C}{p}\right)^{k}.
\end{equation}
\end{theorem}

\subsection{Proof}\label{ssec:fnr-proof}

\begin{proof}[Proof of Theorem~\ref{thm:fnr-rccbf}]
We proceed in several steps.

\medskip\noindent\textbf{Step 1: Count at an inserted element's position.}

Element $x$ contributes exactly one increment to position $h_i(x)$
(through hash function $h_i$) during the Insert operation.
Other elements $x' \in S \setminus \{x\}$ may also hash to the same
position.  Under the Poisson hashing model, the number of
\emph{other} elements hashing to position $h_i(x)$ through any of the $k$
hash functions is $\Poi(kn/m - k/m) = \Poi(\lambda - k/m)$.  For large
$m$, $k/m \to 0$, so the total count at position $h_i(x)$ is
\begin{equation}\label{eq:ci-dist}
  C_i \;\sim\; 1 + \Poi\!\left(\lambda - \frac{k}{m}\right)
  \;\approx\; 1 + \Poi(\lambda).
\end{equation}
In particular, $C_i \ge 1$ always (since $x$ itself contributes at
least~$1$).

Note that the ``$1+$'' accounts for the deterministic contribution of
element~$x$ itself.  Both self-collisions (element $x$ hashing to the same
position via different hash functions, which occurs with probability
$O(k^2/m)$) and cross-element collisions (other elements hashing to the
same position) contribute to the remaining Poisson count.  The total
discrepancy from treating these independently is $O(k/m)$, which we
suppress throughout.

\medskip\noindent\textbf{Step 2: Read probability at a single position.}

Given that position $h_i(x)$ has total count $C_i$, the Read operation
returns $R_i = p \cdot S_j$ for a uniformly random sub-counter $S_j$.
By the same argument as Lemma~\ref{lem:read-threshold}:
\begin{equation}\label{eq:read-xi}
  \Prob(R_i \ge \theta \mid C_i)
  = s_\tau(C_i / p).
\end{equation}
A false negative at position $i$ occurs when $R_i < \theta$, i.e.:
\begin{equation}\label{eq:fn-per-position}
  \Prob(\text{FN at position } i \mid C_i)
  = 1 - s_\tau(C_i/p).
\end{equation}
For $\tau = 1$:
\begin{equation}\label{eq:fn-per-position-tau1}
  \Prob(\text{FN at position } i \mid C_i)
  = 1 - \tanh(C_i/p).
\end{equation}

\medskip\noindent\textbf{Step 3: Independence across positions.}

The reads at positions $h_1(x), \ldots, h_k(x)$ are performed on
\emph{different} MultiCounters (since the hash functions map to different
positions).  Conditional on the counts $(C_1, \ldots, C_k)$, the Read
operations are independent: each selects a uniformly random sub-counter
within its own MultiCounter.  Therefore:
\begin{align}
  \Prob(\text{no FN} \mid C_1, \ldots, C_k)
  &= \prod_{i=1}^{k} \Prob(R_i \ge \theta \mid C_i) \notag \\
  &= \prod_{i=1}^{k} s_\tau(C_i/p).
  \label{eq:no-fn-product}
\end{align}
The false negative rate is one minus this:
\begin{equation}\label{eq:fnr-final}
  \FNR = 1 - \prod_{i=1}^{k} s_\tau(C_i/p).
\end{equation}
This proves~\eqref{eq:fnr-rccbf}.

\medskip\noindent\textbf{Step 4: Specialisation to $\tau = 1$ and equal counts.}

For $\tau = 1$, using $s_1(\mu) = \tanh(\mu)$
(Lemma~\ref{lem:s1-tanh}):
\begin{equation}\label{eq:fnr-tau1-proof}
  \FNR = 1 - \prod_{i=1}^{k}\tanh(C_i/p).
\end{equation}
If $C_i = C$ for all $i$:
\begin{equation}\label{eq:fnr-equal-proof}
  \FNR = 1 - \tanh(C/p)^k.
\end{equation}

\medskip\noindent\textbf{Step 5: Expected FNR.}

The counts $C_i$ are random.  From Step~1, each $C_i$ has distribution
$1 + \Poi(\lambda)$ (approximately).  The expected FNR is
\begin{equation}\label{eq:expected-fnr}
  \E[\FNR]
  = 1 - \E\!\left[\prod_{i=1}^{k}\tanh(C_i/p)\right].
\end{equation}
Under the Poisson model, the counts at distinct positions are independent
(by the same argument as Lemma~\ref{lem:independence}), so:
\begin{equation}\label{eq:expected-fnr-product}
  \E\!\left[\prod_{i=1}^{k}\tanh(C_i/p)\right]
  = \prod_{i=1}^{k} \E[\tanh(C_i/p)]
  = \left(\E[\tanh(C/p)]\right)^k,
\end{equation}
where $C \sim 1 + \Poi(\lambda)$.  Therefore:
\begin{equation}\label{eq:expected-fnr-final}
  \E[\FNR]
  = 1 - \left(\sum_{c=1}^{\infty}
        \frac{e^{-\lambda}\,\lambda^{c-1}}{(c-1)!}\;
        \tanh\!\left(\frac{c}{p}\right)\right)^{k}.
\end{equation}
\end{proof}

\subsection{Behavioural Analysis of the FNR}\label{ssec:fnr-behaviour}

\begin{proposition}[Low-Load FNR]\label{prop:fnr-lowload}
When $C/p \ll 1$ (low effective load per sub-counter):
\begin{equation}\label{eq:fnr-lowload}
  \Prob(\text{FN at one position}) \;\approx\; 1 - \frac{C}{p}.
\end{equation}
In particular, if $C = 1$ (only the element itself hashes to the position)
and $p \gg 1$, the per-position false negative probability approaches~$1$.
\end{proposition}

\begin{proof}
For small $x$, $\tanh(x) = x - x^3/3 + O(x^5) \approx x$.
Setting $x = C/p$ with $C/p \ll 1$:
\begin{equation}\label{eq:tanh-small}
  \tanh(C/p) \approx C/p.
\end{equation}
Therefore the per-position false negative probability is
\begin{equation}\label{eq:fn-small}
  1 - \tanh(C/p) \approx 1 - C/p.
\end{equation}
When $C = 1$ and $p \gg 1$, this is approximately $1 - 1/p \approx 1$,
meaning a false negative is almost certain at each position.
\end{proof}

\begin{proposition}[High-Load FNR]\label{prop:fnr-highload}
When $C/p \gg 1$ (high effective load per sub-counter):
\begin{equation}\label{eq:fnr-highload}
  \FNR \;\to\; 0.
\end{equation}
\end{proposition}

\begin{proof}
For large $x$, $\tanh(x) \to 1$ with $1 - \tanh(x) = 2e^{-2x}/(1 +
e^{-2x}) \le 2e^{-2x}$.  Setting $x = C/p$ with $C/p \gg 1$, we have
$\tanh(C/p) \to 1$ exponentially fast, so $\tanh(C/p)^k \to 1$ and
$\FNR = 1 - \tanh(C/p)^k \to 0$.
\end{proof}

\begin{proposition}[Minimum Load for Bounded FNR]\label{prop:min-load}
For $\FNR \le \delta$ (with equal counts $C_i = C$), it is necessary and
sufficient that
\begin{equation}\label{eq:min-load}
  C \;\ge\; p \cdot \mathrm{arctanh}\!\left((1 - \delta)^{1/k}\right).
\end{equation}
\end{proposition}

\begin{proof}
From~\eqref{eq:fnr-equal}, $\FNR \le \delta$ iff
$\tanh(C/p)^k \ge 1 - \delta$,
iff $\tanh(C/p) \ge (1-\delta)^{1/k}$,
iff $C/p \ge \mathrm{arctanh}((1-\delta)^{1/k})$.
The last step uses the fact that $\tanh$ is strictly increasing on
$(0, \infty)$ and that $(1-\delta)^{1/k} < 1$ for $\delta > 0$ and
$k \ge 1$, ensuring $\mathrm{arctanh}$ is well-defined.
\end{proof}

\begin{remark}[The FPR--FNR Trade-Off]\label{rem:tradeoff}
Propositions~\ref{prop:sparse} and~\ref{prop:fnr-lowload} reveal the
fundamental trade-off:
\begin{itemize}
  \item Increasing $p$ \emph{reduces} FPR (by a factor of $(1/p)^k$ in
        the sparse regime) but \emph{increases} FNR (requiring $C \ge p
        \cdot \mathrm{arctanh}((1-\delta)^{1/k})$ for bounded FNR).
  \item The crossover occurs around $C \approx p$: when the total count at
        a position is comparable to the number of sub-counters, both FPR
        and FNR are moderate.
\end{itemize}
This is the price of relaxation: better false positive performance comes at
the cost of potential false negatives if the load is insufficient.
\end{remark}

%%%%%%%%%%%%%%%%%%%%%%%%%%%%%%%%%%%%%%%%%%%%%%%%%%%%%%%%%%%%
%% SECTION 8 — SPECIAL CASES AND COMPARISON
%%%%%%%%%%%%%%%%%%%%%%%%%%%%%%%%%%%%%%%%%%%%%%%%%%%%%%%%%%%%
\section{Special Cases and Comparison}\label{sec:special}

We verify that the \RCCBF{} analysis reduces to known results in limiting
cases and examine the extremes of the parameter space.

\begin{corollary}[$p = 1$: Recovery of Standard CBF]\label{cor:p1}
When $p = 1$ and $\theta = 1$:
\begin{align}
  \FPR_{\RCCBF} &= (1 - e^{-\lambda})^k = \FPR_{\mathrm{BF}},
  \label{eq:fpr-p1} \\
  \FNR_{\RCCBF} &= 0. \label{eq:fnr-p1}
\end{align}
\end{corollary}

\begin{proof}
When $p = 1$, the MultiCounter has one sub-counter.  Read returns
$R = C$ exactly.  Query returns ``present'' if and only if $C \ge 1$.

\emph{FPR.}  For $y \notin S$, the count at each queried position is
$C \sim \Poi(\lambda)$.  A false positive occurs when all $k$ positions
satisfy $C \ge 1$, so
$\FPR = \Prob(C \ge 1)^k = (1 - e^{-\lambda})^k = \FPR_{\mathrm{BF}}$.

\emph{FNR.}  For $x \in S$, $C_i \ge 1$ at every position (since $x$
itself contributes at least one increment).  Hence $R_i = C_i \ge 1
= \theta$ at every position, and $\FNR = 0$.
\end{proof}

\begin{corollary}[$\theta = 0$: Trivial FPR]\label{cor:theta0}
When $\theta = 0$:
\begin{equation}\label{eq:fpr-theta0}
  \FPR_{\RCCBF} = 1.
\end{equation}
\end{corollary}

\begin{proof}
With $\theta = 0$, the condition $R_i \ge 0$ is always satisfied (since
sub-counters are non-negative).  Every query returns ``present'',
regardless of membership.
\end{proof}

\begin{corollary}[Large $p$, Small $\lambda$]\label{cor:large-p-small-lam}
When $p \to \infty$ with $\lambda = O(1)$:
\begin{align}
  \FPR_{\RCCBF} &\approx \left(\frac{\lambda}{p}\right)^k \to 0
  \qquad\text{(dramatic reduction)}, \label{eq:fpr-large-p} \\
  \FNR_{\RCCBF} &\to 1
  \qquad\text{(catastrophic false negative rate)}.
  \label{eq:fnr-large-p}
\end{align}
\end{corollary}

\begin{proof}
\emph{FPR.}  From Proposition~\ref{prop:sparse}:
$q(\lambda, p, 1) \approx \lambda/p$ for small $\lambda$ and large $p$.
Hence $\FPR \approx (\lambda/p)^k \to 0$.

\emph{FNR.}  For an inserted element $x$ with $C_i \ge 1$ at each
position, $\tanh(C_i/p) \approx C_i/p$.  With $C_i = O(1)$ and
$p \to \infty$:
$\prod_{i=1}^{k}\tanh(C_i/p) \approx \prod_{i=1}^{k}(C_i/p) \to 0$.
Hence $\FNR = 1 - \prod \tanh(C_i/p) \to 1$.
\end{proof}

\begin{corollary}[Large $\lambda/p$: Convergence to Standard
CBF]\label{cor:large-load}
When $\lambda / p \to \infty$, both the FPR and FNR of the \RCCBF{}
converge to those of the standard CBF:
\begin{align}
  \FPR_{\RCCBF} &\to (1 - e^{-\lambda})^k, \label{eq:fpr-converge} \\
  \FNR_{\RCCBF} &\to 0. \label{eq:fnr-converge}
\end{align}
\end{corollary}

\begin{proof}
When $\lambda / p \gg 1$, the Poisson count $C$ satisfies $C/p \gg 1$
with high probability, so $\tanh(C/p) \to 1$.  In the FPR expression,
$q = \E[\tanh(C/p)] \to \Prob(C \ge 1) = 1 - e^{-\lambda}$ (the gap
$1 - \tanh(C/p)$ is exponentially small for $C \ge 1$).  The FNR
$= 1 - \tanh(C/p)^k \to 0$ since each factor tends to~$1$.
The \RCCBF{} thus behaves as a standard CBF in the heavily loaded
regime.
\end{proof}

\begin{remark}[Summary of Regimes]\label{rem:regimes}
Table~\ref{tab:regimes} summarises the behaviour across parameter regimes.

\begin{table}[h]
\centering
\begin{tabular}{lcc}
\hline
\textbf{Regime} & \textbf{FPR vs.\ Standard CBF} & \textbf{FNR} \\
\hline
$p = 1$ (standard CBF) & Equal & $0$ \\
$\lambda \ll 1$, $p \gg 1$ & $\approx (1/p)^k \times$ smaller
  & $\approx 1 - (1/p)^k$ \\
$\lambda \approx p$ (crossover) & Moderately smaller & Moderate \\
$\lambda \gg p$ & $\to$ Equal & $\to 0$ \\
$\theta = 0$ & FPR $= 1$ & FNR $= 0$ \\
\hline
\end{tabular}
\caption{FPR and FNR behaviour of the \RCCBF{} across parameter regimes.}
\label{tab:regimes}
\end{table}
\end{remark}

%%%%%%%%%%%%%%%%%%%%%%%%%%%%%%%%%%%%%%%%%%%%%%%%%%%%%%%%%%%%
%% SECTION 9 — CONCLUSION
%%%%%%%%%%%%%%%%%%%%%%%%%%%%%%%%%%%%%%%%%%%%%%%%%%%%%%%%%%%%
\section{Conclusion}\label{sec:conclusion}

We have presented a rigorous analysis of the false positive and false
negative rates of the Relaxed Concurrent Counting Bloom Filter (\RCCBF{}),
a data structure that distributes each counter across $p$ sub-counters
updated via the power-of-two-choices paradigm.

\paragraph{Main results.}
In the mean-field limit ($p \to \infty$), the FPR is
$q(\lambda, p, \tau)^k$ (Theorem~\ref{thm:fpr-rccbf}), where the
per-position false positive probability $q$ is expressed as a Poisson
average of the occupancy tail fractions $s_\tau$.  For the threshold-one
case, we obtained a convergent closed-form series
(Theorem~\ref{thm:fpr-tau1}) and proved that the \RCCBF{} strictly
improves upon the standard Bloom filter FPR for all $p > 1$
(Theorem~\ref{thm:fpr-strict}).  The FNR is
$1 - \prod_{i=1}^{k} s_\tau(C_i/p)$ (Theorem~\ref{thm:fnr-rccbf}).

\paragraph{Key trade-off.}
Increasing the relaxation parameter $p$ reduces the FPR by a factor of
approximately $(1/p)^k$ in the sparse regime but introduces a non-zero FNR
that approaches~$1$ when the effective load $C/p$ is small.  The crossover
regime $\lambda \approx p$ offers a balanced operating point.

\paragraph{Limitations.}
Our analysis relies on three simplifying assumptions:
\begin{enumerate}
  \item The mean-field limit $p \to \infty$, with finite-$p$ convergence
        rates only characterised to $O(1/\sqrt{p})$ (Remark~\ref{rem:finite-p-rate}).
  \item The insertion-only model; we do not analyse the effect of
        deletions on MultiCounter occupancy or FPR/FNR.
  \item Sequential allocation; the mean-field ODE does not account for
        stale reads under concurrent execution
        (Remark~\ref{rem:concurrent-validity}).
\end{enumerate}

\paragraph{Future directions.}
Natural extensions include:
(i)~tight finite-$p$ error bounds via Stein's method or coupling arguments,
(ii)~analysis of deletions and their impact on the occupancy distribution,
(iii)~a concurrent model incorporating stale-read probabilities and their
effect on load balance, and
(iv)~experimental validation on multi-core hardware to compare theoretical
predictions with observed FPR, FNR, and throughput.

%%%%%%%%%%%%%%%%%%%%%%%%%%%%%%%%%%%%%%%%%%%%%%%%%%%%%%%%%%%%
%% BIBLIOGRAPHY
%%%%%%%%%%%%%%%%%%%%%%%%%%%%%%%%%%%%%%%%%%%%%%%%%%%%%%%%%%%%
\begin{thebibliography}{10}

\bibitem{bloom1970}
B.~H.~Bloom,
\newblock Space/time trade-offs in hash coding with allowable errors,
\newblock \emph{Communications of the ACM}, 13(7):422--426, 1970.

\bibitem{azar1999}
Y.~Azar, A.~Z.~Broder, A.~R.~Karlin, and E.~Upfal,
\newblock Balanced allocations,
\newblock \emph{SIAM Journal on Computing}, 29(1):180--200, 1999.

\bibitem{mitzenmacher1996}
M.~Mitzenmacher,
\newblock \emph{The Power of Two Choices in Randomized Load Balancing},
\newblock PhD thesis, University of California, Berkeley, 1996.

\bibitem{mitzenmacher2001}
M.~Mitzenmacher,
\newblock The power of two choices in randomized load balancing,
\newblock \emph{IEEE Transactions on Parallel and Distributed Systems},
  12(10):1094--1104, 2001.

\bibitem{mitzenmacher2001survey}
M.~Mitzenmacher, A.~Richa, and R.~Sitaraman,
\newblock The power of two random choices: A survey of techniques and results,
\newblock In \emph{Handbook of Randomized Computing}, pages 255--312,
  Kluwer, 2001.

\bibitem{kurtz1970}
T.~G.~Kurtz,
\newblock Solutions of ordinary differential equations as limits of pure jump
  Markov processes,
\newblock \emph{Journal of Applied Probability}, 7(1):49--58, 1970.

\bibitem{peres2015}
Y.~Peres, K.~Talwar, and U.~Wieder,
\newblock Graphical balanced allocations and the $(1+\beta)$-choice process,
\newblock \emph{Random Structures \& Algorithms}, 47(4):760--775, 2015.

\bibitem{alistarh2018}
D.~Alistarh, J.~Aspnes, M.~Gelashvili, and N.~Shavit,
\newblock Distributionally linearizable data structures,
\newblock In \emph{Proceedings of the 30th ACM Symposium on Parallelism in
  Algorithms and Architectures (SPAA)}, pages 133--142, 2018.

\bibitem{fan2000}
L.~Fan, P.~Cao, J.~Almeida, and A.~Z.~Broder,
\newblock Summary cache: A scalable wide-area web cache sharing protocol,
\newblock \emph{IEEE/ACM Transactions on Networking}, 8(3):281--293, 2000.

\bibitem{berenbrink2006}
P.~Berenbrink, A.~Czumaj, A.~Steger, and B.~V{\"o}cking,
\newblock Balanced allocations: The heavily loaded case,
\newblock \emph{SIAM Journal on Computing}, 35(6):1350--1385, 2006.

\end{thebibliography}

\end{document}
